\clearpage\section{问题一算法的复杂度与优化}\label{appendix-a}

\subsection{算法的计算量}

每次测向产生两个半平面约束，\(m\) 次观测共给出 \(2m\) 条边界。算法1枚举边界对，候选交点数量至多为 \(\binom{2m}{2}\)；每个交点需要代入全部 \(2m\) 个约束进行检验，因此顶点枚举阶段的计算量为 \(O(m^3)\)。该估计针对边界求交与约束校验，不包含可行性及有界性判别所调用的线性规划求解器的内部计算。

得到 \(h\) 个顶点后，枚举顶点对计算直径的计算量为 \(O(h^2)\)，检验直径圆是否覆盖全部顶点的计算量为 \(O(h)\)。对于由 \(2m\) 个半平面构成的有界凸多边形，\(h\le2m\)，因此这两个步骤不改变顶点枚举部分的立方级计算量。

\subsection{检测点较多时的优化}

当检测点数量较多时，可将边界对枚举替换为按方向排序的半平面求交。先将半平面按边界方向排序，再用双端队列维护有效边界：同向平行约束保留较严格者，新约束加入时从队列两端删除与之不相容的边界。排序需要 \(O(m\log m)\) 的计算量，每条边界至多入队和出队一次，后续维护为线性计算量；空集、无界及退化情形仍须分别判别。

对求得的有界凸多边形，可进一步采用旋转卡壳法，沿有序顶点遍历对踵点对，在 \(O(h)\) 的计算量内求出直径，再按正文判据检验覆盖性。这样，半平面求交与直径计算这两个主要几何步骤可达到 \(O(m\log m)\) 的总计算复杂度。

正文算例使用的是边界对枚举实现。本节所述排序求交与旋转卡壳属于可选优化方案，尚未用于所报告的实验，因此不据此作出实际运行时间已经缩短的结论。
